\documentclass[12pt,a4paper]{article}

\usepackage[utf8]{inputenc}
\usepackage[T1]{fontenc}
\usepackage[margin=2.5cm]{geometry}
\usepackage{amsmath,amssymb}
\usepackage{booktabs}
\usepackage{array}
\usepackage{hyperref}
\usepackage{microtype}

\hypersetup{
  colorlinks=true,
  linkcolor=blue!60!black,
  urlcolor=blue!60!black,
}

\title{\textbf{Datathon 2026 --- Ensemble Forecasting of Revenue and COGS}}
\author{}
\date{April 2026}

\begin{document}

\maketitle

\begin{abstract}
We present an expanding-window cross-validated ensemble of three classical time-series forecasters
(Theta, Holt-Winters, Hybrid ARIMA) to forecast daily Revenue and COGS for 2023--2024. Ensemble
weights are jointly optimized across three chronological folds (2020, 2021, 2022) to minimize a
weighted objective of RMSE (45\%), MAE (35\%), and $(1 - R^2)$ (20\%). Classical statistical
models are chosen over tree-based or neural approaches because the test horizon requires
extrapolation beyond the training distribution (2012--2022). Average cross-validated performance:
MAE = 522,663, RMSE = 731,302, $R^2 = 0.7778$. All code, hyperparameters, and reproducibility
details are provided.
\end{abstract}

\section{Method: Pipeline, Validation, and Leakage Control}

\subsection{Model Ensemble and Feature Engineering}

The ensemble combines three classical time-series forecasters:
\begin{itemize}
  \item \textbf{Theta}: Decomposition method with log scaling and learned day-of-week adjustment.
  \item \textbf{Holt-Winters}: Triple exponential smoothing with multiplicative seasonality ($m = 7$ weeks).
  \item \textbf{Hybrid ARIMA}: Linear seasonal baseline (calendar + geometric growth) plus ARIMA$(p,0,q)$ on recent residuals (last 3 years).
\end{itemize}

ARIMA orders are pre-tuned via grid search on expanding-window CV: Revenue ARIMA$(2,0,3)$,
COGS ARIMA$(2,0,3)$. All models use cyclic encoding of calendar features:
\[
  \text{doy\_sin} = \sin\left(\frac{2\pi \cdot \text{doy}}{365.25}\right), \quad
  \text{doy\_cos} = \cos\left(\frac{2\pi \cdot \text{doy}}{365.25}\right),
\]
and similarly for day-of-week with period 7, to preserve cyclic continuity at year/week boundaries.

Post-ensemble, predictions are scaled, adjusted for day-of-week effects, biased, clipped to
$\geq 0$, capped at COGS $\leq 1.05 \times$ Revenue, and rounded to 2 decimal places.
All post-processing parameters (scales, biases) are calibrated on CV folds only.

\subsection{Expanding-Window Cross-Validation}

Unlike $k$-fold CV, expanding-window CV preserves temporal order: training window grows with
each fold, validation year never seen during training. This prevents leakage and preserves
long-run growth trends critical for extrapolation.

\begin{table}[h!]
\centering
\small
\begin{tabular}{lrr|rrr}
\toprule
\textbf{Fold} & \textbf{Train end} & \textbf{Val period} & \textbf{MAE} & \textbf{RMSE} & \textbf{$R^2$} \\
\midrule
2020 & 2019-12-31 & 2020 & 520,078 & 707,297 & 0.7845 \\
2021 & 2020-12-31 & 2021 & 507,854 & 741,795 & 0.7703 \\
2022 & 2021-12-31 & 2022 & 540,058 & 744,813 & 0.7787 \\
\midrule
\textbf{Avg} & — & — & \textbf{522,663} & \textbf{731,302} & \textbf{0.7778} \\
\bottomrule
\end{tabular}
\caption{Expanding-window CV results. Metrics computed on concatenated $[\text{Revenue}, \text{COGS}]$ vector.}
\end{table}

\subsection{Leakage Control and Ensemble Weights}

To prevent test-set contamination:
(1) Each fold trains strictly on $\text{Date} < \text{fold\_year}$; test Revenue/COGS never loaded.
(2) All statistics (scale factors, day-of-week profiles, interpolation) computed per-fold.
(3) Ensemble weights optimized on out-of-fold validation errors, not the test horizon.
(4) Sample submission CSV used only for future dates; its Revenue/COGS columns ignored.

Weights minimized via Nelder-Mead with 12 random restarts over softmax-parameterized vectors:
\begin{table}[h!]
\centering
\small
\begin{tabular}{lcc}
\toprule
\textbf{Model} & \textbf{Revenue} & \textbf{COGS} \\
\midrule
Theta & 0.1742 & 0.1670 \\
Holt-Winters & 0.0000 & 0.0000 \\
Hybrid ARIMA & 0.8258 & 0.8330 \\
\bottomrule
\end{tabular}
\caption{Optimized ensemble weights from expanding-window CV.}
\end{table}

\section{Submission and Reproducibility}

Final forecast: 548 days (2023-01-01 to 2024-07-01).
Total Revenue: 2.322B VND, Total COGS: 2.109B VND, COGS/Revenue: 0.908.

All code available in \texttt{scripts/forecast.py} and \texttt{src/} modules.
Run \texttt{python scripts/forecast.py} to regenerate submission and this report.

\newpage
\appendix

\section{Detailed Feature Engineering and Cyclic Encoding}

\subsection{Input Preprocessing}

Daily sales data sorted chronologically. Missing values within training window linearly interpolated,
then forward/back-filled, applied per-fold independently. All values clipped to non-negative.

\subsection{Per-Model Features}

\begin{table}[h!]
\centering
\small
\begin{tabular}{llll}
\toprule
\textbf{Model} & \textbf{Feature} & \textbf{Encoding} & \textbf{Purpose} \\
\midrule
Theta & $\log(1 + y)$ & Log transform & Variance stabilization \\
Theta & DoW (0--6) & Multiplicative factor & Weekly rhythm \\
HW & Raw $y$ & ETS fit & Level, trend, seasonal \\
ARIMA & (month, day) & Calendar mean profile & Intra-year baseline \\
ARIMA & Year index & Geometric growth & Trend extrapolation \\
ARIMA & 3-yr residuals & ARIMA fit & Autocorrelation \\
\bottomrule
\end{tabular}
\caption{Feature engineering strategy per model.}
\end{table}

\subsection{Cyclic Encoding Rationale}

Raw integer features create artificial boundaries: doy 365 and doy 1 treated as maximally
distant despite being one day apart. Cyclic $(\sin, \cos)$ encoding preserves distance:
\[
  \text{dow\_sin} = \sin\left(\frac{2\pi \cdot \text{dow}}{7}\right), \quad
  \text{dow\_cos} = \cos\left(\frac{2\pi \cdot \text{dow}}{7}\right).
\]

One-hot encoding avoided: would create 365 or 7 sparse columns and still fail to encode cyclic
distance correctly. $(\sin, \cos)$ pair is minimal sufficient representation.

\section{Complete Leakage Control Table}

\begin{table}[h!]
\centering
\small
\begin{tabular}{p{3.5cm}p{7cm}}
\toprule
\textbf{Control Point} & \textbf{Implementation} \\
\midrule
No future Revenue/COGS in training &
  Each fold filters strictly to $\text{Date} < \text{fold\_year}$; test Revenue/COGS never
  loaded at any stage. \\
\addlinespace
No global statistics from test data &
  Scale factors, means, and standard deviations computed only from each fold's training
  portion. \\
\addlinespace
DoW profiles recomputed per fold &
  DoW multipliers use only $\text{train}[\text{Date.year} < \text{fold\_year}]$; fold 2022
  profile does not see 2022 actuals. \\
\addlinespace
Interpolation within fold only &
  Missing values filled using each fold's training window alone; no forward-fill from future
  observations. \\
\addlinespace
Ensemble weights from out-of-fold validation &
  Softmax weights optimized on fold validation errors, not on 2023--2024 test horizon. \\
\addlinespace
Sample submission used only for dates &
  Test CSV supplies future dates only; its Revenue/COGS columns are never used. \\
\bottomrule
\end{tabular}
\caption{Six leakage control measures applied at every pipeline stage.}
\end{table}

\section{Feature Importance and Explainability}

\subsection{Why not SHAP?}

SHAP (TreeSHAP, DeepSHAP, KernelSHAP) requires specific model classes or treats time steps
as independent samples, discarding temporal autocorrelation. Classical forecasters (Theta, ARIMA)
exploit temporal structure explicitly; Pearson correlation is deterministic and interpretable
alternative.

\subsection{Importance Scores}

Absolute Pearson correlation with concatenated Revenue + COGS:

\begin{table}[h!]
\centering
\small
\begin{tabular}{lrl}
\toprule
\textbf{Feature} & \textbf{Corr (\%)} & \textbf{Interpretation} \\
\midrule
doy\_cos & 16.86\% & Cyclic annual phase (peak at seasonal high) \\
is\_month\_end & 13.17\% & End-of-month purchasing spike (payroll cycle) \\
days\_since & 10.91\% & Long-term growth trend (largest driver of level) \\
doy\_sin & 10.68\% & Cyclic annual phase (rising Jan$\to$Jul) \\
year & 10.31\% & Annual level shift (year-over-year growth) \\
dom & 9.81\% & Day-of-month position \\
month & 6.87\% & Broad seasonal periods (summer/winter) \\
week & 6.36\% & Within-year seasonal progression \\
\bottomrule
\end{tabular}
\caption{Calendar feature importance (absolute Pearson correlation).}
\end{table}

Key insights: (1) Cyclic annual features dominate due to strong annual cycle in fashion
e-commerce. (2) End-of-month effect ($\approx 13\%$) indicates consistent purchasing spikes
correlated with payroll cycles. (3) Trend component ($\approx 21\%$ combined) is largest driver
of absolute revenue level — why tree models fail at extrapolation. (4) Monthly periodicity
reinforced across multiple features.

\section{Mathematical Formulas}

\subsection{Competition Metrics}

\[
  \text{MAE} = \frac{1}{n} \sum_{i=1}^{n} |F_i - A_i|, \quad
  \text{RMSE} = \sqrt{\frac{1}{n} \sum_{i=1}^{n} (F_i - A_i)^2},
\]
\[
  R^2 = 1 - \frac{\sum_{i=1}^{n}(A_i - F_i)^2}{\sum_{i=1}^{n}(A_i - \bar{A})^2},
\]
where $n = 2 \times 365 = 730$ (Revenue and COGS concatenated) and $\bar{A}$ is mean of actuals.

\subsection{Theta Method}

Decomposes $y_t$ into modified theta lines:
\[
  \tilde{y}_0(t) = 2\bar{y} - y_t \quad (\text{no seasonality, retains trend}),
  \quad \tilde{y}_2(t) = y_t \quad (\text{full series}).
\]
Forecast blends simple exponential smoothing on $\tilde{y}_0$ with OLS linear trend:
\[
  F(h) = \tfrac{1}{2}\, \text{SES}(\tilde{y}_0, h) + \tfrac{1}{2}(a + b(T + h)).
\]
Applied on $\log(1 + y)$; output exponentiated and multiplied by learned day-of-week factor.

\subsection{Holt-Winters (Multiplicative Seasonality)}

\begin{align*}
  L_t &= \alpha \frac{y_t}{S_{t-m}} + (1-\alpha)(L_{t-1} + B_{t-1}), \\
  B_t &= \beta (L_t - L_{t-1}) + (1-\beta) B_{t-1}, \\
  S_t &= \gamma \frac{y_t}{L_t} + (1-\gamma) S_{t-m}, \\
  F(h) &= \left(L_T + h B_T\right) \, S_{T+h - m\lceil h/m \rceil},
\end{align*}
with $m = 7$ (weekly seasonality) and $\alpha, \beta, \gamma$ optimized to minimize sum of
squared errors.

\subsection{Hybrid ARIMA}

\textbf{Step 1 --- Seasonal Baseline:}
\[
  \hat{y}(t) = b \cdot g^{\Delta_{\text{yr}}} \cdot s(\text{month}, \text{day}),
\]
where $b$ = mean daily level (last training year), $g$ = geometric mean growth rate (last 3 years),
and $s$ = seasonal profile $\bar{y}_{(\text{month,day})} / \bar{y}_{\text{annual}}$.

\textbf{Step 2 --- Residual ARIMA$(p,0,q)$:}
\[
  e_t = y_t - \hat{y}(t), \quad
  e_t = c + \sum_{i=1}^{p} \varphi_i e_{t-i} + \sum_{j=1}^{q} \theta_j \varepsilon_{t-j} + \varepsilon_t.
\]
Fit on residuals from last 3 training years only.

\textbf{Step 3 --- Combine:}
\[
  F(h) = \max\left(0, \hat{y}(h) + \text{ARIMA\_forecast}(h)\right).
\]

\subsection{Ensemble Objective}

Let $\mathbf{w}_{\text{rev}} = \text{softmax}(\mathbf{z}_{1:3})$ and
$\mathbf{w}_{\text{cog}} = \text{softmax}(\mathbf{z}_{4:6})$ for unconstrained
$\mathbf{z} \in \mathbb{R}^6$. For each fold year $k$:
\[
  \mathcal{L}_k = 0.45 \cdot \frac{\text{RMSE}_k}{735{,}000}
                + 0.35 \cdot \frac{\text{MAE}_k}{532{,}000}
                + 0.20 \cdot (1 - R^2_k).
\]
\[
  \min_{\mathbf{z}}\; \frac{1}{3}\sum_{k \in \{2020,2021,2022\}} \mathcal{L}_k
  \quad \text{via Nelder-Mead, 12 random restarts}.
\]

\section{Complete Hyperparameter Table}

\begin{table}[h!]
\centering
\small
\begin{tabular}{lll}
\toprule
\textbf{Parameter} & \textbf{Value} & \textbf{Tuning Method} \\
\midrule
ARIMA Revenue order $(p, d, q)$ & $(2, 0, 3)$ & Grid search on CV folds \\
ARIMA COGS order $(p, d, q)$ & $(2, 0, 3)$ & Grid search on CV folds \\
ARIMA residual window (years) & 3 & CV \\
Holt-Winters trend & Additive & Fixed \\
Holt-Winters seasonality & Multiplicative & Fixed \\
Holt-Winters seasonal\_periods & 7 & Fixed (weekly) \\
Theta period & 365 & Fixed (annual) \\
Theta log transform & True & Fixed \\
SCALE\_REVENUE & 1.184 & CV calibration \\
SCALE\_COGS & 1.191 & CV calibration \\
DOW\_STRENGTH & 0.40 & Manual tuning \\
REVENUE\_BIAS (VND/day) & 25,000 & CV calibration \\
COGS\_BIAS (VND/day) & 112,500 & CV calibration \\
MAX\_COGS\_RATIO & 1.05 & Competition constraint \\
Optimizer & Nelder-Mead & Fixed \\
Optimizer restarts & 12 & Fixed \\
Random seed & 42 & Fixed \\
\bottomrule
\end{tabular}
\caption{Complete hyperparameter table with tuning provenance.}
\end{table}

\section{Final Results and Model Comparison}

\subsection{Submission Totals}

\begin{table}[h!]
\centering
\small
\begin{tabular}{lrr}
\toprule
\textbf{Metric} & \textbf{Revenue} & \textbf{COGS} \\
\midrule
Total 2023--2024 & 2.322B VND & 2.109B VND \\
Daily average & 4,238,128 VND & 3,848,980 VND \\
COGS / Revenue & — & 0.908 \\
\bottomrule
\end{tabular}
\caption{Final submission aggregates for 2023--2024 test period (548 days).}
\end{table}

\subsection{Model Comparison and Selection Rationale}

Classical models chosen because tree-based and neural models systematically fail on extrapolation.
LightGBM and random forests cap predictions within training range; a growing series in test
period is under-predicted. Chronos-T5 zero-shot failed on 2-year horizon. Classical methods
with explicit trend components (geometric growth, exponential smoothing) extrapolate reliably
to unseen future states.

\begin{thebibliography}{9}
\bibitem{Assimakopoulos2000}
  V.~Assimakopoulos and K.~Nikolopoulos,
  ``The theta model: a decomposition approach to forecasting,''
  \emph{International Journal of Forecasting}, vol.~16, no.~4, pp.~521--530, 2000.

\bibitem{Hyndman2021}
  R.~Hyndman and G.~Athanasopoulos,
  ``Forecasting: principles and practice,'' 3rd ed.,
  OTexts, Melbourne, Australia, 2021.
  Available at \url{https://otexts.com/fpp3/}.
\end{thebibliography}

\end{document}
